% Pursuing Stacks (cotes 134-2 to 134-8): archive cote 134-7, site page 101 (PDF page 102), typescript page 543.
% typescript. Draft compared: ps6.tex:5092-5148.
% AI-transcribed by Claude Opus 5.5 (claude-opus-5-5, adaptive thinking, effort high), Message Batches API, batch ps4 (msgbatch_012Aoc1NuskATPsyBeKSYK7J), 2026-11-12; not yet checked by eye.
\scanpage{134-7/101}
\mspage{543}

to introduce a notion of ``q\xout{h}uasi-coherent \inserted{homological} quasi-isomorphism'' between
such objects, as a map (17) such that the corresponding map (17')
should be a quasi-isomorphism, which presumably may be interpreted also
in cohomological terms\hand{,} as usual. As just noticed, it is doubtful that
this is implied by (17) being a weak equivalence, and even if it should
be implied, it looks considerably weaker in a way - just as in the
set-up of usual homotopy types, \xout{it} a homology equivalence is a conside-%
rably weaker notion than homotopy equivalence, unless \xout{the} we make \xout{the}
a 1-connectedness assumption. We could reinforce of course this notion
of \inserted{(quasi-coherent)} ``homological'' or ``cohomological'' quasi-isomorphism, by taking to the
cohomological version of it, and instead of quasi-coherent \xout{complexes of}
coefficients \uline{coming from the base} S=Spec(k), admit equally the analog
of ``\uline{twisted coefficients}'', which would amount here to taking as coeffi-%
cients quasi-coherent sheaves \inserted{$F_*$} on $Y_*$ , such that for any structural
morphism
\[ \phi : Y_m \to Y_n \]
associated to a map in \hand{$\Delta$} , the corresponding map of \xout{sheaves} quasi-cohe-%
rent sheaves on $Y_m$
\[ \phi^*(F_n) \to F_m \]
should be an isomorphism. (These``twisted coefficients''should correspond
to quasi-coherent sheaves on S=Spec(k), \uline{on which the sheaf of groups
$\underline{\Pi}_1(Y_*)$ ope\unsure{ra}tes}, in case at least when $Y_*$ is relatively 0-connected
and endowed with a section, so that $\underline{\Pi}_1(Y_*)$ makes a sense.) My point
here is that it may be interesting to take the derived category of a
suitable category of semisimplicial schemes (submitted to some \inserted{very mild} conditi-%
ons, such as quasi-compactness and quasi-separation say) with respect
\margin{q.c.h. (circled)} to this notion of q\hand{.}c\hand{.}h\hand{.} quasi-isomorphism - with the hope that the set of
maps with respect to which we are now localizing is wide enough, so as
to get \uline{the same} derived category as when working \inserted{only} with ``bundles'', namely
getting just schematic homotopy types. This would give a very strong
link between more or less arbitary semisimplicial schemes over k (and
stronger still when k=$\underline{Z}$)\hand{,} and ordinary homotopy types, of a much subtler
nature than the known link with ordinary pro-homotopy types via etale
cohomology with discrete coefficients. But maybe \xout{this is} the daydrea-%
ming is getting here so much out of reach or maybe simply crazy, that
I better stop along this line !

\margin{101}
